\section{Locally ringed spaces}
\label{section-locally-ringed-spaces}

\noindent
Recall that we defined ringed spaces
in Sheaves, Section \ref{sheaves-section-ringed-spaces}.
Briefly, a ringed space is a pair $(X, \mathcal{O}_X)$ consisting
of a topological space $X$ and a sheaf of rings $\mathcal{O}_X$.
A morphism of ringed spaces $f : (X, \mathcal{O}_X) \to (Y, \mathcal{O}_Y)$
is given by a continuous map $f : X \to Y$ and an $f$-map of sheaves
of rings $f^\sharp : \mathcal{O}_Y \to \mathcal{O}_X$. You can
think of $f^\sharp$ as a map $\mathcal{O}_Y \to f_*\mathcal{O}_X$, see
Sheaves, Definition \ref{sheaves-definition-f-map} and
Lemma \ref{sheaves-lemma-f-map}.

\medskip\noindent
A good geometric example of this to keep in mind is
$\mathcal{C}^\infty$-manifolds and morphisms of
$\mathcal{C}^\infty$-manifolds. Namely, if $M$ is a
$\mathcal{C}^\infty$-manifold, then the sheaf $\mathcal{C}^\infty_M$
of smooth functions is a sheaf of rings on $M$. And any
map $f : M \to N$ of manifolds is smooth if and only if
for every local section $h$ of $\mathcal{C}^\infty_N$
the composition $h \circ f$ is a local section of $\mathcal{C}^\infty_M$.
Thus a smooth map $f$ gives rise in a natural way to a morphism
of ringed spaces
$$
f : (M , \mathcal{C}^\infty_M) \longrightarrow (N, \mathcal{C}^\infty_N)
$$
see Sheaves, Example \ref{sheaves-example-continuous-map-ringed}.
It is instructive to consider what happens to stalks. Namely, let
$m \in M$ with image $f(m) = n \in N$. Recall that the stalk
$\mathcal{C}^\infty_{M, m}$ is the ring of germs of smooth
functions at $m$, see
Sheaves, Example \ref{sheaves-example-germs-functions}.
The algebra of germs of functions on $(M, m)$ is a local ring with maximal
ideal the functions which vanish at $m$.
Similarly for $\mathcal{C}^\infty_{N, n}$. The map on stalks
$f^\sharp : \mathcal{C}^\infty_{N, n} \to \mathcal{C}^\infty_{M, m}$
maps the maximal ideal into the maximal ideal, simply
because $f(m) = n$.

\medskip\noindent
In algebraic geometry we study schemes. On a scheme the sheaf of rings
is not determined by an intrinsic property of the space.
The spectrum of a ring $R$
(see Algebra, Section \ref{algebra-section-spectrum-ring}) endowed
with a sheaf of rings constructed out of $R$ (see below),
will be our basic building block. It will turn out that
the stalks of $\mathcal{O}$ on $\Spec(R)$ are the local
rings of $R$ at its primes. There are two reasons to introduce
locally ringed spaces in this setting: (1) There is in general
no mechanism that assigns to a continuous map of spectra a
map of the corresponding rings. This is why we add as an extra datum
the map $f^\sharp$. (2) If we consider morphisms
of these spectra in the category of ringed spaces, then the
maps on stalks may not be local homomorphisms. Since our
geometric intuition says it should we introduce locally ringed spaces
as follows.

\begin{definition}
\label{definition-locally-ringed-space}
Locally ringed spaces.
\begin{enumerate}
\item A {\it locally ringed space $(X, \mathcal{O}_X)$}
is a pair consisting of a
topological space $X$ and a sheaf of rings $\mathcal{O}_X$ all of whose stalks
are local rings.
\item Given a locally ringed space $(X, \mathcal{O}_X)$ we say that
$\mathcal{O}_{X, x}$ is the {\it local ring of $X$ at $x$}.
We denote $\mathfrak{m}_{X, x}$ or simply $\mathfrak{m}_x$
the maximal ideal of $\mathcal{O}_{X, x}$. Moreover, the
{\it residue field of $X$ at $x$} is the residue field
$\kappa(x) = \mathcal{O}_{X, x}/\mathfrak{m}_x$.
\item A {\it morphism of locally ringed spaces}
$(f, f^\sharp) : (X, \mathcal{O}_X) \to (Y, \mathcal{O}_Y)$
is a morphism of ringed spaces such that for all $x \in X$
the induced ring map $\mathcal{O}_{Y, f(x)} \to \mathcal{O}_{X, x}$ is a
local ring map.
\end{enumerate}
\end{definition}

\noindent
We will usually suppress the sheaf of rings $\mathcal{O}_X$
in the notation when discussing locally ringed spaces. We will
simply refer to ``the locally ringed space $X$''.
We will by abuse of notation think of $X$ also as
the underlying topological space. Finally we will denote
the corresponding sheaf of rings
$\mathcal{O}_X$ as the {\it structure sheaf of $X$}.
In addition, it is customary to denote the maximal ideal
of the local ring $\mathcal{O}_{X, x}$ by
$\mathfrak{m}_{X, x}$ or simply $\mathfrak{m}_x$.
We will say ``let $f : X \to Y$ be a morphism of locally ringed
spaces'' thereby suppressing the structure sheaves even further.
In this case, we will by abuse of notation think of $f : X\to Y$
also as the underlying continuous map of topological spaces.
The $f$-map corresponding to $f$ will customarily
be denoted $f^\sharp$. The condition that $f$ is a morphism
of locally ringed spaces can then be expressed by saying that
for every $x\in X$ the map on stalks
$$
f^\sharp_x : \mathcal{O}_{Y, f(x)} \longrightarrow \mathcal{O}_{X, x}
$$
maps the maximal ideal $\mathfrak m_{Y, f(x)}$ into
$\mathfrak m_{X, x}$.

\medskip\noindent
Let us use these notational conventions to show that the
collection of locally ringed spaces and morphisms of locally
ringed spaces forms a category. In order to see this we have
to show that the composition of morphisms of locally ringed
spaces is a morphism of locally ringed spaces. OK, so let
$f : X \to Y$ and $g : Y \to Z$ be morphism of locally
ringed spaces. The composition of $f$ and $g$ is defined in
Sheaves, Definition \ref{sheaves-definition-composition-maps-ringed-spaces}.
Let $x \in X$. By
Sheaves, Lemma \ref{sheaves-lemma-compose-f-maps-stalks}
the composition
$$
\mathcal{O}_{Z, g(f(x))}
\xrightarrow{g^\sharp}
\mathcal{O}_{Y, f(x)}
\xrightarrow{f^\sharp}
\mathcal{O}_{X, x}
$$
is the associated map on stalks for the morphism $g \circ f$.
The result follows since a composition of local ring
homomorphisms is a local ring homomorphism.

\medskip\noindent
A pleasing feature of the definition is the fact that the functor
$$
\textit{Locally ringed spaces}
\longrightarrow
\textit{Ringed spaces}
$$
reflects isomorphisms (plus more).
Here is a less abstract statement.

\begin{lemma}
\label{lemma-isomorphism-locally-ringed}
\begin{slogan}
An isomorphism of ringed spaces between locally ringed spaces is an
isomorphism of locally ringed spaces.
\end{slogan}
Let $X$, $Y$ be locally ringed spaces.
If $f : X \to Y$ is an isomorphism of
ringed spaces, then $f$ is an isomorphism
of locally ringed spaces.
\end{lemma}

\begin{proof}
This follows trivially from the corresponding fact in algebra:
Suppose $A$, $B$ are local rings. Any isomorphism of rings
$A \to B$ is a local ring isomorphism.
\end{proof}













